% ======================================================================
% ABSTRACT
% Final prose. Approximately 900 characters.
% ======================================================================
\noindent
On 28 April 2026, the U.S. Food and Drug Administration announced two
real-time clinical trial proofs-of-concept and a pilot Request for
Information, framing the conduct of trials in which key safety signals
no longer take years to reach the agency. The
prevailing oncology AI patient-prediction baseline remains a
heterogeneous set of narrow supervised models with ceilings such as
Manz 2020 AUC 0.89, the SHIELD-RT prospective randomized trial,
SCORPIO, and PROGPATH, set against the Huang 2025 null result that
machine learning provides no significant gain over Cox regression on
real-world structured survival data. This paper uses four author
Physical AI oncology trial simulations - Simulation 1 in
\path{hour-00} through \path{hour-55}, Simulation 2 ten patient-journey
stages, Simulation 3 a 24-hour autonomous sponsor, and Simulation 4 a
168-hour 7-day sponsor extension verified locally on a Core i5-6200U
laptop with 4 GB RAM - to demonstrate that Claude Code Opus 4.7 Max
produces working agentic code applicable to patient prediction better than supervised models in current trial practice and with more utility than the FDA
RTCT proof-of-concept. The computational
signature - 1M token contextual code-and-text awareness, hourly commit
cadence, repository-scale forecasting - is the source of the advantage,
without claiming clinical deployment readiness.
